\section*{第 \ref{ch:b2:equilibrium-shifts} 章——平衡常数与平衡移动}

\begin{solution}{exo:b2:equilibrium-shifts:1}
$K^\circ = \exp(32\,800/(8.314 \times 298.15)) = \num{5.6e5}$，而对于
$+\qty{32.8}{kJ/mol}$ 的反应，其逆反应为 \num{1.8e-6}。若使平衡常数增大为 100 倍，
则需要 $\Delta(\Delta_r G^\circ) = -RT\ln 100 = \qty{-11.4}{kJ/mol}$。
\end{solution}

\begin{solution}{exo:b2:equilibrium-shifts:2}
$\ln K^\circ(400) = \ln(5.4 \times 10^5) + (91\,880/8.314)(1/400 -
1/298.15) = 13.20 - 9.44 = 3.76$，$K^\circ \approx 43$。查表得 35.6：
在 \qty{100}{K} 范围内，视 $\Delta_r H^\circ$ 为常数的近似已经
偏差 20\,\%（随 $T$ 升高，反应变得更放热）。
\end{solution}

\begin{solution}{exo:b2:equilibrium-shifts:3}
(a) $n = 1$，$r = 0$，$\varphi = 2$：$v = 1$（蒸气压是
$T$ 的函数）。(b) $v = 3 - 1 + 2 - 3 = 1$。(c) $v = 3 - 1 + 2 - 1 = 3$。(d) 四种
物种，一个反应，两个相（碳与气体）：$v = 4 - 1 + 2 - 2 =
3$；若气体仅由碳与水蒸气生成，$n(\ce{CO}) = n(\ce{H2})$，
在气相中，$s = 1$，$v = 2$。
\end{solution}

\begin{solution}{exo:b2:equilibrium-shifts:4}
加热：吸热方向，即向 \ce{NO2} 一侧（正向）。压缩：
朝气体分子数减少的方向，即 \ce{N2O4} 一侧（逆向）。恒容加入氩气：
平衡不移动，反应气体的分压不变。
\end{solution}

\begin{solution}{exo:b2:equilibrium-shifts:5}
在 \qty{298}{K} 时：$\Delta_r G^\circ = 180.6 - 298.15 \times 0.0248 =
\qty{173.2}{kJ/mol}$，$K^\circ = \num{4.5e-31}$。在 \qty{2000}{K} 时：$\Delta_r
G^\circ = 180.6 - 2000 \times 0.0248 = \qty{131.0}{kJ/mol}$, $K^\circ =
\num{3.8e-4}$。当 $\Delta\nu_{\text{气体}} = 0$ 时，$x(\ce{NO})^2 = K^\circ x(\ce{N2})
x(\ce{O2})$: $x(\ce{NO}) = \sqrt{3.8 \times 10^{-4} \times 0.78 \times 0.21} =
\num{7.9e-3}$，几乎为百分之一。热机在如此高温下运行；气体冷却时，
$K^\circ$ 急剧减小，但 \ce{NO} 的分解
十分缓慢：高温平衡被冻结在废气中（直至催化剂
将其去除）。
\end{solution}

\begin{solution}{exo:b2:equilibrium-shifts:6}
若酯为 $\xi$ 摩尔、总物质的量恒定，则 $Q =
\xi^2/[(1 - \xi)(n_B - \xi)]$。当 $n_B = 1.0$ 时：$\xi^2/(1 - \xi)^2 = 4$，$\xi
= \qty{0.67}{mol}$。当 $n_B = 3.0$ 时：$3\xi^2 - 16\xi + 12 = 0$，$\xi =
\qty{0.90}{mol}$。移去水使 $Q$ 始终低于 $K^\circ$，无论 $\xi$ 为何值，
反应持续正向进行，直至酸耗尽。
\end{solution}

\begin{solution}{exo:b2:equilibrium-shifts:7}
无氩气时：物质的量分别为 $1 - \alpha$、$2\alpha$，总物质的量为 $1 + \alpha$；
$Q = 4\alpha^2/(1 - \alpha^2) = 0.148$, $\alpha = \sqrt{0.148/4.148} = 0.189$.
加入 \qty{1}{mol} 氩气，压强为 \qty{1}{bar}，总物质的量为 $2 + \alpha$，$Q =
4\alpha^2/[(1 - \alpha)(2 + \alpha)] = 0.148$，故 $4.148\alpha^2 + 0.148\alpha
- 0.296 = 0$，$\alpha = 0.250$：恒总压稀释会降低
分压，有利于气体分子数较多的一侧。在
恒容下，反应气体的分压不变：
$\alpha$ 仍为 0.189。
\end{solution}

\begin{solution}{exo:b2:equilibrium-shifts:8}
只有固体：$n = 3$，$r = 1$，$s = 1$（$p(\ce{NH3}) = p(\ce{HCl})$），
$\varphi = 2$：$v = 3 - 1 - 1 + 2 - 2 = 1$；选定 $T$ 后，压强
即被确定（称为“解离压”）。加入氨后，$s = 0$，$v = 2$：
可任选 $T$ 与一个压强；乘积 $p(\ce{NH3})\,p(\ce{HCl})$ 是
由 $T$ 决定。
\end{solution}

\begin{solution}{exo:b2:equilibrium-shifts:9}
$\Delta_r H^\circ = -R\ln(K_2/K_1)/(1/T_2 - 1/T_1) = -8.314 \times
\ln(0.0173)/(1/600 - 1/500) = \qty{-101}{kJ/mol}$，比
\qty{298}{K} 时（\qty{-91.9}{kJ/mol}）更负，正如基尔霍夫定律所预言
（$\Delta_r C_p^\circ < 0$）。
\end{solution}

\begin{solution}{exo:b2:equilibrium-shifts:10}
$Q = \dfrac{n_{\ce{NH3}}^2\,n_{\text{总}}^2}{n_{\ce{N2}}n_{\ce{H2}}^3}
\left(\dfrac{p^\circ}{p}\right)^2$。在 $T$、$p$ 恒定时：
$\partial\ln Q/\partial n_{\ce{N2}} = -1/n_{\ce{N2}} + 2/n_{\text{总}}$，为正，
当 $n_{\ce{N2}} > n_{\text{总}}/2$ 时，即 $x(\ce{N2}) > 1/2$。此时 $Q$
超过 $K^\circ$，平衡逆向移动。恒容时，
$\partial\ln Q/\partial n_{\ce{N2}} = -1/n_{\ce{N2}} < 0$：总是正向。
\end{solution}

\begin{solution}{exo:b2:equilibrium-shifts:11}
$K^\circ = \exp(-1620/(8.314 \times 1100)) = 0.84$: $p(\ce{CO2}) =
\qty{0.84}{bar}$，平衡时即 $n = pV/RT = 0.84 \times 10^5 \times
0.0100/(8.314 \times 1100) = \qty{0.092}{mol}$ 的气体。若加入 \qty{0.050}{mol}
石灰石，全部分解（$p = \qty{0.46}{bar} < K^\circ p^\circ$）：
不存在平衡，即平衡被打破。若加入 \qty{0.200}{mol}，则达到平衡：
\qty{0.092}{mol} 的 \ce{CO2} 与 \ce{CaO}，剩余 \qty{0.108}{mol} 的 \ce{CaCO3}。
\end{solution}

\begin{solution}{exo:b2:equilibrium-shifts:12}
第 1 床层出口温度约 \qty{890}{K}，转化率 68\,\%；第 2 床层出口温度约
\qty{785}{K}，转化率 92\,\%；第 3 床层出口温度约 \qty{725}{K}，转化率 98\,\%。在
一个绝热床层中，气体边反应边升温，直至与平衡
曲线相遇，该温度下仅允许约 68\,\% 的转化。冷催化剂
全程使用不可行，因为催化剂只在高温下起作用；段间
冷却既保持反应速率，又使平衡极限上升。
末段床层受限于催化剂仍能工作的最低温度所对应的
平衡（此后工厂吸收三氧化硫，并使气体再次通过
催化剂）。
\end{solution}

\begin{solution}{pb:b2:equilibrium-shifts:1}
\textbf{1.} $\Delta_r H^\circ = \qty{-91.88}{kJ/mol}$, $\Delta_r S^\circ =
2(192.77) - 191.61 - 3(130.68) = \qty{-198.1}{J/(K.mol)}$, $\Delta_r G^\circ =
\qty{-32.8}{kJ/mol}$, $K^\circ = \num{5.6e5}$.
\textbf{2.} $K^\circ(700) = \exp(-54\,380/(8.314 \times 700)) = \num{8.8e-5}$;
$K^\circ(800) = \exp(-77\,320/(8.314 \times 800)) = \num{8.9e-6}$.
\textbf{3.} $\Delta_r H^\circ = -R\ln(K_{800}/K_{700})/(1/800 - 1/700)$,
即 \qty{-106}{kJ/mol}，比 \qty{298}{K} 时更负。
\textbf{4.} $\ln K^\circ(723) = \ln(8.75 \times 10^{-5}) + 12\,777 \times
(1/723.15 - 1/700)$，其中 $12\,777 = 106\,230/8.314$：$K^\circ =
\num{4.9e-5}$.
\textbf{5.} $\Delta_r G^\circ(723) = -91.88 + 723.15 \times 0.1981 =
\qty{51.4}{kJ/mol}$，$K^\circ = \num{1.9e-4}$：偏大了四倍，因为
$\Delta_r H^\circ$ 与 $\Delta_r S^\circ$ 在 \qty{425}{K} 范围内都会改变
($\Delta_r C_p^\circ \approx \qty{-44}{J/(K.mol)}$).
\textbf{6.} \ce{N2} $1 - \xi$，\ce{H2} $3 - 3\xi$，\ce{NH3} $2\xi$，总物质的量 $4 -
2\xi$.
\textbf{7.} $x = 2\xi/(4 - 2\xi)$，故 $1 - x = (4 - 4\xi)/(4 - 2\xi)$；
氮的摩尔分数为 $(1 - \xi)/(4 - 2\xi) = (1 - x)/4$，氢的摩尔分数
是其三倍。
\textbf{8.}
\[
  K^\circ = \frac{x^2}{\frac{1-x}{4}\left(\frac{3(1-x)}{4}\right)^3}
  \left(\frac{p^\circ}{p}\right)^2 = \frac{256}{27}\,\frac{x^2}{(1 - x)^4}
  \left(\frac{p^\circ}{p}\right)^2 ;
\]
两边开平方。
\textbf{9.} $a = (\sqrt{27}/16) \times \sqrt{4.9 \times 10^{-5}} \times 1 =
\num{2.27e-3}$；由 $x/(1 - x)^2 = a$ 得 $x = 0.0023$。
\textbf{10.} $a = 0.455$; $ax^2 - (2a + 1)x + a = 0$: $x = 0.25$.
\textbf{11.} $n = 3$，$r = 1$，$s = 1$（反应保持的比例为 1 : 3），
$\varphi = 1$：$v = 2$。选定 $T$ 与 $p$ 后，便已无自由度。
\textbf{12.} 加压：正向（$\Delta\nu_{\text{气体}} = -2$）。升温：
逆向（放热）。
\textbf{13.} 总压恒定时加入惰性气体，会降低反应气体的
分压，与降低压强的效果相同：逆向。
\textbf{14.} $\partial\ln Q/\partial n_{\ce{N2}} = (-1/0.30 + 2)/n_{\text{总}} < 0$：
$Q$ 低于 $K^\circ$，正向。
\textbf{15.} 否：恒容下 $\partial\ln Q/\partial n_{\ce{N2}} =
-1/n_{\ce{N2}} < 0$，无论组成如何，正向。
\textbf{16.} 氨几乎为零时，入口处 $Q$ 远低于 $K^\circ$：
气体远离平衡进入，并正向反应。
\textbf{17.} 气体在催化剂上停留的时间过短，以致
达不到平衡；延长接触时间则需要更大、更
昂贵的转化器，却只换来少量收益。
\textbf{18.} 由 $x = 2\xi/(4 - 2\xi) = 0.18$ 得 $\xi = \qty{0.305}{mol}$，每
摩尔进料的 \ce{N2}：氢（及氮）的单程转化率为
$0.305$, 31\,\%.
\textbf{19.} 稳态运行时，新鲜进料最终全部转化：
每单位时间有 \qty{2.00}{mol} 氨离开；转化器须通过 $1.00/
0.305 = \qty{3.3}{mol}$ 的 \ce{N2}（每单位时间），其余部分循环。
\textbf{20.} 随进料进入的氩（及甲烷）从不反应，并
不随液氨移除：若不驰放，它们将
在循环中积累，并降低反应物的分压。
\textbf{21.} 在 \qty{500}{K} 时，铁催化剂太慢：
平衡虽有利，却永远无法在工业可接受的时间内接近。
\textbf{22.} 在 \qty{723}{K} 与 \qty{200}{bar} 下，按化学计量比进料，理想
气体：$\boldsymbol{x(\ce{NH3}) \approx 0.25}$。
\end{solution}
